%% front/notice.tex -- attribution + legal + purpose. Hand-written, stable.
\chapter*{Sources, Attribution and Purpose}
\hbtopicmark{SOURCES AND ATTRIBUTION}

\section*{What this book is}

This handbook is an \emph{original synthesis}. It is not a republication,
abridgement, paraphrase-in-place, or substitute for any source it draws on.
Every idea taken from a source is credited to that source by short ID tag
(\texttt{B1}, \texttt{B2}, \dots) at entry level and in the dossier appendix,
and all prose here was written afresh for this handbook. Chapter and law
titles from the sources are used only as short identifying labels.

\section*{Primary sources ingested so far}

\begingroup\small
\begin{itemize}
\item[\textbf{B1}] Sparkman, R.\,B. \emph{The Art of Manipulation}.
  Cope Allied Publishing / Dial Press, 1978.
  Contributed: street-level tactic set (Tactics 1--14), character typologies,
  the financial-dealings chapter, the ``unargue'' method, ingratitude handling.
\item[\textbf{B2}] \emph{The Art of Manipulation} -- condensed study summary
  of B1 (Bookey edition, 2025).
  Contributed: chapter-level outline of B1 used only to cross-check coverage
  of B1. Treated as a \emph{secondary} witness: where B1 and B2 disagree,
  B1 wins. B2 never originates a claim.
\item[\textbf{B3}] Greene, R. \emph{The 48 Laws of Power}.
  Viking/Penguin, 1998.
  Contributed: the 48-law strategic layer, transgression/observance case
  structure, the ``reversal'' device, courtiership and timing doctrine.
\end{itemize}
\par\endgroup

\vspace{4pt}
The authoritative machine-readable list is \texttt{registry/books.yaml}.
Adding a source is documented in \texttt{docs/INGESTION.md}; it is an
append-only operation.

\section*{Why the material is presented this way}

The sources describe techniques that are used on ordinary people every day,
in sales, in offices, in families and in politics. A technique you cannot
name is a technique you cannot refuse. Each entry therefore carries its
\textbf{Counters} slot before its \textbf{Limits} slot, and the recognition
material (\textbf{Tells}) is given as much weight as the execution material
(\textbf{Moves}). The handbook is written to be usable defensively first.

Nothing here is advice to harm, defraud, coerce or exploit anyone.
Several entries describe conduct that is illegal in many jurisdictions
(fraud, extortion, blackmail, misrepresentation). Those entries exist so the
reader can recognise the pattern being used against them; the \textbf{Limits}
slot records the legal and practical blowback. Use this book on yourself and
in your own defence.

\section*{Editorial policy on ``padding''}

The sources spend most of their length on historical anecdote, case
narrative, repetition and moral hedging. All of that is stripped. What is
kept is the transferable mechanism: what the move is, why it lands, how it
shows up, how it is answered, and where it fails. Where a source's only
contribution to an entry is an anecdote, the anecdote is dropped and the
entry is not padded to compensate. See \texttt{docs/STYLE-GUIDE.md}.
